\documentclass[11pt]{book}

\usepackage{hyperref}
\usepackage{listings}
\usepackage{graphicx}
\usepackage{minted}
\usepackage{geometry}
\usepackage{fancyhdr}

\title{PACTS: Permission, Accountability, Compliance, Transparency System}
\author{Project Documentation}
\date{\today}

\begin{document}

\maketitle
\tableofcontents

\chapter{System Overview}
\section{Introduction}
PACTS (Permission, Accountability, Compliance, Transparency System) is an advanced framework for managing AI system safety, compliance, and transparency in software projects.

\section{Core Components}
\begin{itemize}
  \item Permission Management
  \item Accountability Framework
  \item Compliance Engine
  \item Transparency Layer
\end{itemize}

\chapter{Architecture}
\section{System Design}
The system is designed with modularity and extensibility in mind:
\begin{itemize}
  \item Core System
  \item Protocol Layer
  \item Verification Engine
  \item Memory Management
\end{itemize}

\section{Bootstrap Mechanism}
PACTS uses a unique bootstrap mechanism where newer versions are developed and tested using previous stable versions.

\chapter{Integration Guide}
\section{Using PACTS}
Example code for system initialization:
\begin{minted}{typescript}
import { PACTS } from '@pacts/core';

const pacts = new PACTS({
  mode: 'strict',
  logging: true,
  compliance: ['AI_SAFETY', 'GDPR']
});
\end{minted}

\section{Version Management}
Guidelines for managing PACTS versions in your project...

\chapter{Development Guide}
\section{Setting Up Development Environment}
Steps for setting up a development environment...

\section{Testing}
Guidelines for testing PACTS components...

\chapter{API Reference}
Detailed API documentation...

\end{document} 